\begin{helper}
For the canonical pre-terminal replay transcript family over \(I\), the
recorded coordinate set attached to each representative is exactly
\(\noderef{TwoBitePreTerminalRecordedPairs}\) for that representative.  More
precisely, there is a finite transcript family with cells \(C_i\), recorded
sets \(R_i\), and representatives \(\omega_i\), satisfying the covering,
disjointness, representative-membership, cylinder, and event-containment
properties of \(\noderef{FixedSetPreTerminalReplayTranscriptConstruction}\),
and also
\[
e\in \noderef{TwoBitePreTerminalRecordedPairs}(\omega_i,\varepsilon,I)
\quad\Longleftrightarrow\quad
e\in R_i
\]
for every transcript index \(i\) and every terminal coordinate \(e\).
\end{helper}
\begin{proof}
Construct the transcript family by the deterministic pre-terminal replay of
\(\noderef{FixedSetPreTerminalReplayTranscriptConstruction}\).  Its first
five conclusions are exactly the construction conclusions from that node:
the cells cover the projection-size event, are pairwise disjoint, contain
their representatives, are described by common injection and common truth
table on \(R_i\), and remain inside the projection-size event.

It remains to identify the recorded set of a representative.  Fix a
representative \(\omega_i\).  The replay inserts a terminal coordinate into
\(R_i\) exactly when one of the pre-terminal queries reads that coordinate.
During the \(F_2\)-pass, the same-color query from a projection-contained
endpoint to an \(F_1\)-endpoint is a pre-terminal recorded coordinate by
\(\noderef{FixedSetPreTerminalF2QueriesRecorded}\).  After the staged-revealed
set is known, every incident query from a staged-revealed endpoint to a
projection-contained endpoint is precisely the incident branch in
\(\noderef{TwoBitePreTerminalRecordedPairs}\).  Finally, each later query of
a pair inside \(X_x(I)\) for a staged-revealed base vertex \(x\) is precisely
the \(X_x(I)\)-pair branch in
\(\noderef{TwoBitePreTerminalRecordedPairs}\), with the red projection in the
red-tagged case and the blue projection in the blue-tagged case.

Conversely, every coordinate satisfying one of the two branches of
\(\noderef{TwoBitePreTerminalRecordedPairs}(\omega_i,\varepsilon,I)\) is
queried by this deterministic replay.  If it satisfies the incident branch,
the endpoint that is staged-revealed has already entered the replay's staged
set, so the ordered incident-coordinate scan reads that terminal coordinate.
If it satisfies the \(X_x(I)\)-pair branch, then the replay's final
\(X_x(I)\)-pair scan for the staged-revealed vertex \(x\) reads it.  The
outer terminal-coordinate condition is exactly the replay's restriction to
valid red- and blue-tagged terminal coordinates.  Hence \(R_i\) and
\(\noderef{TwoBitePreTerminalRecordedPairs}(\omega_i,\varepsilon,I)\) have
the same elements.
\end{proof}
